\part{Symmetry and Quantum Numbers}

\chapter{Symmetry separates state space into sectors}
\label{ch:representations}

\physicalquestion{Why do energy levels occur in symmetry-related families
rather than as unrelated individual states?}

\modelassumptions{A symmetry is an invertible transformation preserving the
physical predictions of the model.  The transformations compose as a finite
group $G$ and act linearly on a finite-dimensional complex state space $V$.}

\section{How does a physical transformation become algebra?}

A group representation is a homomorphism
$\rho:G\to\operatorname{GL}(V)$.  A subspace $W\subseteq V$ is invariant
when $\rho(g)W\subseteq W$ for all $g$.  It is irreducible when its only
invariant subspaces are $0$ and itself.  A map $T:V\to W$ intertwines two
representations when $T\rho_V(g)=\rho_W(g)T$.

If a Hamiltonian $H$ obeys $H\rho(g)=\rho(g)H$, then every eigenspace of $H$
is invariant.  Symmetry can therefore organize degeneracy only if invariant
spaces can be decomposed into elementary pieces.

\section{Why can finite symmetry be completely decomposed?}

\begin{theorem}[Maschke and Schur]
Every finite-dimensional complex representation of a finite group is a direct
sum of irreducible representations.  If $V,W$ are irreducible, every nonzero
intertwiner $V\to W$ is an isomorphism; moreover
$\End_G(V)=\C\,1_V$.
\end{theorem}

\begin{proof}
Average any Hermitian inner product over $G$ to make it invariant.  If
$W\subseteq V$ is invariant, then $W^\perp$ is invariant as well.  Induction
on dimension gives a decomposition into irreducibles.  For an intertwiner
$T$, its kernel and image are invariant, so irreducibility makes a nonzero
$T$ invertible.  If $T\in\End_G(V)$, choose an eigenvalue $\lambda$ over
$\C$.  The intertwiner $T-\lambda1$ has nonzero kernel and is therefore zero.
\end{proof}

\section{What does decomposition predict physically?}

Write
\[
V\cong\bigoplus_\alpha M_\alpha\otimes V_\alpha,
\]
where the $V_\alpha$ are inequivalent irreducibles and $G$ acts trivially on
the multiplicity spaces $M_\alpha$.  Every symmetry-preserving Hamiltonian
has the form $\bigoplus_\alpha H_\alpha\otimes1_{V_\alpha}$.  Thus an energy
inside the $\alpha$ sector has at least $\dim V_\alpha$-fold degeneracy unless
other structure splits the multiplicity space.  The claim is conditional on
the symmetry being exact; perturbations that break $G$ may split the level.

\section{What further contact should be proved?}

\begin{exercise}[Averaging produces invariant projectors]
For $W\subseteq V$ invariant and any projection $P:V\to W$, prove that
\[
\widetilde P=\frac1{|G|}\sum_{g\in G}\rho(g)P\rho(g)^{-1}
\]
is an equivariant projection onto $W$.
\end{exercise}

\begin{exercise}[Symmetry protects degeneracy]
Let an irreducible representation $V_\alpha$ occur once in an eigenspace of a
$G$-invariant Hamiltonian.  Prove that every $G$-invariant perturbation acts
there by a scalar.  Explain why this conclusion says nothing about a
perturbation that breaks $G$.
\end{exercise}


\chapter{Characters predict molecular spectra}
\label{ch:characters}

\physicalquestion{Why are some molecular vibrations visible in infrared or
Raman spectra while symmetry forbids others?}

\modelassumptions{The equilibrium molecule has a finite point group $G$.
Small nuclear displacements form a linear representation, and the electric
dipole approximation is valid for the transitions under discussion.}

\section{Which invariant detects an irreducible sector?}

The character of a representation is
$\chi_V(g)=\tr(\rho_V(g))$.  It is constant on conjugacy classes and does not
depend on a basis.  For class functions define
\[
\langle\chi,\psi\rangle_G
=\frac1{|G|}\sum_{g\in G}\overline{\chi(g)}\psi(g).
\]

For an irreducible $V_\alpha$ of dimension $d_\alpha$, the operator
\[
P_\alpha=\frac{d_\alpha}{|G|}
\sum_{g\in G}\overline{\chi_\alpha(g)}\rho_V(g)
\]
is the candidate projector onto its isotypic component.

\section{Why do characters recover multiplicities?}

\begin{theorem}[Character orthogonality]
Irreducible characters of a finite group are orthonormal.  If
$V\cong\bigoplus_\alpha m_\alpha V_\alpha$, then
\[
m_\alpha=\langle\chi_\alpha,\chi_V\rangle_G,
\]
and $P_\alpha$ is the projection onto the $V_\alpha$-isotypic component.
\end{theorem}

\begin{proof}
For a linear map $A:V_\alpha\to V_\beta$, average
$A\mapsto |G|^{-1}\sum_g\rho_\beta(g)A\rho_\alpha(g)^{-1}$.
The result is an intertwiner.  Schur's lemma from
Chapter~\ref{ch:representations} makes it zero for $\alpha\ne\beta$ and a
scalar identity for $\alpha=\beta$.  Applying this to matrix units and taking
traces gives character orthogonality.  Additivity of trace on direct sums
then gives the multiplicity formula.  Substitution into $P_\alpha$ shows that
it is identity on copies of $V_\alpha$ and zero on all other irreducibles.
\end{proof}

\section{How does this return to a water spectrum?}

For water at equilibrium the model point group is $C_{2v}$.  The nuclear
displacement representation has dimension nine.  Removing translations and
rotations leaves
\[
V_{\mathrm{vib}}\cong2A_1\oplus B_2.
\]
A dipole transition from $V_i$ to $V_f$ through a dipole component $V_\mu$
can be nonzero only if
\[
\Hom_G(1,V_f^*\otimes V_\mu\otimes V_i)\ne0.
\]
This invariant-tensor condition predicts allowed lines.  Measured line
positions additionally depend on masses, force constants, anharmonicity, and
instrumental conditions; symmetry organizes visibility but does not determine
the frequencies by itself.  Reference vibrational frequencies provide the
experimental comparison for this classification \cite{shimanouchi}.

\section{What further contact should be proved?}

\begin{exercise}[Water has three symmetry-classified vibrations]
Construct the $C_{2v}$ character of the nine-dimensional displacement space
from the fixed nuclei and spatial traces.  Subtract translations and rotations
and derive $2A_1\oplus B_2$.
\end{exercise}

\begin{exercise}[Raman selection differs from dipole selection]
Let the polarizability transform as $\operatorname{Sym}^2(V_\mu)$.  Determine
which $C_{2v}$ irreducibles occur and compare the Raman-active and
infrared-active water modes.
\end{exercise}


\chapter{Continuous symmetry produces angular momentum}
\label{ch:lie}

\physicalquestion{Why are rotational quantum numbers discrete even though
physical rotations vary continuously?}

\modelassumptions{Rotations act continuously and unitarily on a
finite-dimensional state space.  Infinitesimal generators are defined on the
whole space; domain issues of unbounded operators are postponed.}

\section{How is continuous transformation differentiated?}

A matrix Lie group is a closed subgroup of $\operatorname{GL}(n,\C)$.  Its
Lie algebra consists of the matrices $X$ for which $e^{tX}$ remains in the
group, with bracket $[X,Y]=XY-YX$.  Differentiating a smooth representation
$U:G\to U(H)$ gives a Lie-algebra representation.  For rotations, define
self-adjoint generators $J_i$ by differentiating the three one-parameter
rotation subgroups; with $\hbar=1$ they satisfy
\[
[J_i,J_j]=i\varepsilon_{ijk}J_k.
\]
Set $J^2=J_1^2+J_2^2+J_3^2$.

\section{Which operator labels rotational sectors?}

\begin{theorem}[The rotational Casimir]
The operator $J^2$ commutes with every $J_i$ and hence with the connected
rotation representation.  On every irreducible complex rotation sector it is
a scalar.
\end{theorem}

\begin{proof}
Using $[AB,C]=A[B,C]+[A,C]B$,
\[
[J^2,J_\ell]
=i\sum_{i,k}\varepsilon_{i\ell k}(J_iJ_k+J_kJ_i)=0,
\]
because the coefficient is antisymmetric in $i,k$ while the parenthesis is
symmetric.  Thus $J^2$ intertwines the Lie-algebra action and therefore the
connected group action.  Schur's lemma makes it scalar on an irreducible
sector.
\end{proof}

\section{Why do angular-momentum quantum numbers result?}

Introduce $J_\pm=J_1\pm iJ_2$.  The commutators
$[J_3,J_\pm]=\pm J_\pm$ make $J_\pm$ raising and lowering operators.  In a
finite-dimensional unitary irreducible sector, positivity of
$J_\mp J_\pm$ forces a highest weight $j\in\tfrac12\Z_{\ge0}$ and weights
$m=-j,-j+1,\ldots,j$, while $J^2=j(j+1)$.  Integer $j$ descends to an honest
representation of $SO(3)$; half-integer $j$ is an honest representation of
its double cover $SU(2)$ and only a projective representation of $SO(3)$.
The continuous group therefore produces discrete labels through unitarity and
finite dimensionality.

\section{What further contact should be proved?}

\begin{exercise}[Spin one-half]
Set $J_i=\sigma_i/2$ for the Pauli matrices.  Verify the commutation
relations, compute $J^2$, and prove that a $2\pi$ rotation acts by $-1$ while
a $4\pi$ rotation acts by $1$.
\end{exercise}

\begin{exercise}[Rotational invariance conserves angular momentum]
Let canonical position and momentum operators satisfy
$[Q_i,P_j]=i\delta_{ij}$ and set
$L_i=\sum_{jk}\varepsilon_{ijk}Q_jP_k$.  For
$H=(P_1^2+P_2^2+P_3^2)/(2m)+V(Q_1^2+Q_2^2+Q_3^2)$, prove
$[H,L_i]=0$ and identify precisely where rotational invariance of the
potential is used.
\end{exercise}


\chapter{Quantum phases require projective symmetry}
\label{ch:projective}

\physicalquestion{How can a symmetry act consistently on physical rays even
when its operators fail to multiply exactly as the group does?}

\modelassumptions{State vectors differing by a phase represent the same pure
state.  A discrete symmetry group acts unitarily on rays, and phases belong to
$U(1)$ with trivial group action.}

\section{Which phases are removable?}

A projective representation has operators satisfying
\[
U_gU_h=\alpha(g,h)U_{gh},\qquad \alpha(g,h)\in U(1).
\]
Associativity forces the cocycle equation
\[
\alpha(g,h)\alpha(gh,k)=\alpha(h,k)\alpha(g,hk).
\]
Rephasing $U_g$ by $\lambda(g)$ changes $\alpha$ by the coboundary
$\lambda(g)\lambda(h)\lambda(gh)^{-1}$.  The quotient of cocycles by
coboundaries is $H^2(G,U(1))$.

\section{Why does cohomology classify projective actions?}

\begin{theorem}[Projective classification]
Equivalence classes of normalized factor systems for projective unitary
representations of $G$ are classified by $H^2(G,U(1))$.  Tensor product
multiplies the corresponding cohomology classes.
\end{theorem}

\begin{proof}
The displayed associativity calculation says exactly that a factor system is
a normalized $2$-cocycle.  Changing representatives of rays rephases the
operators and hence multiplies the cocycle by a coboundary; conversely every
coboundary is produced by such a rephasing.  For tensor products,
$(U_g\otimes V_g)(U_h\otimes V_h)=
\alpha(g,h)\beta(g,h)U_{gh}\otimes V_{gh}$, so classes multiply.
\end{proof}

\section{How can a cohomology class protect an edge state?}

Let $G=\Z/2\times\Z/2$ be generated by two $\pi$ rotations.  On a spin-half
edge take $U_x=\sigma_1$ and $U_z=\sigma_3$.  Their anticommutation realizes
the nontrivial element of $H^2(G,U(1))\cong\Z/2$.  Any local edge operator
commuting with both is scalar, so a symmetry-preserving local perturbation
cannot split the doublet.  In the spin-one Haldane model, the statement relies
on a bulk gap, locality, and preservation of the specified symmetry; closing
the gap or breaking the symmetry evades the conclusion.

\section{What further contact should be proved?}

\begin{exercise}[The projective class squares to zero]
Compute the factor system for $U_x,U_z$ above and prove that its tensor square
is cohomologically trivial.  Interpret this as the possibility of combining
two protected edges into an ordinary linear representation.
\end{exercise}

\begin{exercise}[Part II synthesis]
Let a Hamiltonian commute with a finite symmetry group and suppose one
isotypic sector carries a projective edge action.  Use characters to isolate
the sector, Schur's lemma to constrain symmetric perturbations, and
$H^2(G,U(1))$ to decide whether the action can be rephased to an ordinary
representation.  Separate degeneracy forced by an irreducible dimension from
degeneracy protected by a nontrivial projective class.
\end{exercise}

\parttransition{Global symmetry uses the same group at every point.  A field
must compare locally varying values and frames across spacetime.  Making
symmetry local forces differential forms, metrics, bundles, connections, and
curvature.}
